%
% Reduktion des Integrals
%
\subsection{Reduktion des Integrals}
Mit dem Werkzeug der quadratfreien Faktorisierung lässen sich jetzt
die für den rationalen Teil des Integrals nötigen Schritte des Plans
vom Ende von Abschnitt~\ref{buch:ratint:section:bernoulli} durchführen.
Zu diesem Zwck muss erst eine Partialbruchzerlegung gefunden werden.
Anschliessend müssen die einzelnen Terme der Partialbruchzerlegung
integriert werden.

%
% Partialbruchzerlegung zur quadratfreien Faktorisierung
%
\subsubsection{Partialbruchzerlegung zur quadratfreien Faktorisierung}
Der wesentliche Schritt zur Konstruktion der Partialbruchzerlegung
ist der folgende Satz.

\begin{satz}
Sei 
\[
f(z)
=
\frac{p(z)}{q(z)^m r(z)},
\]
wobei $q(z)$ teilerfremd zu $r(z)$ ist.
Dann gibt es 
\[
f(z)
=
f_0(z)
+
\frac{s(z)}{r(z)}
+
\sum_{k=1}^m \frac{t_k(z)}{q(z)^k},
\]
wobei $f_0(z)$ ein Polynom ist, $s(z)$ ein Polynom mit Grad
$\deg s(z)<\deg r(z)$ und die Polynome $t_k(z)$ Grad $\deg t_k(z)<\deg q(z)$
haben.
\end{satz}

\begin{proof}
Da $q(z)$ und $r(z)$ teilerfremd sind, sind auch $q(z)^m$ und $r(z)$
teilerfremd.
Daher gibt es $s(z)$ und $t(z)$ mit
\begin{equation}
p(z)
=
s(z)q(z)^m + t(z)r(z).
\label{buch:ratint:hermite:eqn:p=sq+tr}
\end{equation}
Einsetzen in den Bruch ergibt
\begin{align}
f(z)
&=
\frac{p(z)}{q(z)^m r(z)}
=
\frac{s(z)q(z)^m + t(z)r(z)}{q(z)^mr(z)}
=
\frac{s(z)}{r(z)}
+
\frac{t(z)}{q(z)^m}.
\label{buch:ratint:hermite:eqn:srtq}
\end{align}
Falls der Zähler einen Grad mindestens so gross wie der Nenner
hat, kann durch Division des Zählers mit Rest das Polynom $f_0(z)$
abgespaltet werden.

Es muss jetzt nur noch gezeigt werden, dass der zweite Bruch in
\eqref{buch:ratint:hermite:eqn:srtq}
in eine Summe der verlangten Art umgeformt werden kann.
Dazu schreibt man den Zähler mithilfe der Polynomdivision als
\begin{equation*}
t(z)
=
a_m(z) q(z) + t_{m}(z)
\end{equation*}
mit $\deg b(z) < \deg q(z)$.
Dann wird
\begin{align*}
\frac{t(z)}{q(z)^m}
&=
\frac{a(z)q(z)+b(z)}{q(z)^m}
=
\frac{a(z)}{q(z)^{m-1}}
+
\frac{b(z)}{q(z)^m}.
\intertext{Durch wiederholte Anwendung dieser Beobachtung auf
den neuen Zähler $a_m(z) = a_{m-1}(z)q(z)+t_{m-1}(z)$
wird daraus eine Summe der Form}
&=
\frac{t_1(z)}{q(z)}
+
\frac{t_2(z)}{q(z)^2}
+
\ldots
+
\frac{t_m(z)}{q(z)^m}.
\end{align*}
Damit ist der Satz vollständig bewiesen.
\end{proof}

Durch Anwendung auf eine Faktorisierung von $r(z)$ kann schrittweise
die Partialbruchzerlegung mit den Teilnennern $q_l(z)^k$, $k\le l$
für alle $l$ gefunden werden.
Man beachte, dass dazu wieder nur die Polynomdivision
und für die Identität 
\eqref{buch:ratint:hermite:eqn:p=sq+tr}
der euklidische Algorithmus benötigt werden.
Beide können für Polynome in einer beliebigen Körpererweiterung
für den Koeffizientenkörper durchgeführt werden.

\begin{beispiel}
\label{buch:ratint:hermite:bsp:partialbruch}
Man finde die Partialbruchzerlegung von
\[
\frac{r(z)}{q(z)}
\quad\text{mit}\quad
\left\{\;
\renewcommand{\arraycolsep}{2pt}
\begin{array}{rcl}
r(z) &=& 6651z^4 + 3969z^2 - 37338z + 21854 \\
q(z) &=& 9z^6 + 6z^5 - 65z^4 + 20z^3 + 135z^2 - 154z + 49.
\end{array}\right.
\]
Man verwendet dafür die quadratfreie Faktorisierung
\[
q(z)
=
(z-1)^4(3z+7)^2,
\]
die in Beispiel~\ref{buch:ratint:hermite:bsp:qf2} gefunden worden war.
Zunächst sind die beiden Faktoren $(z-1)^4$ und $(3z+7)^2$
teilerfremd, so dass es Polynome $s(z)$ und $t(z)$ gibt derart,
dass
\begin{align*}
r(z) &= s(z)(z-1)^4 + t(z) (3z+7)^2
\intertext{nämlich}
s(z) &=
	\frac{1}{50000}(3214890^5 z + 10180485 z^4 + 1928934 z^3 - 12037977 z^2
	- 46842138 z\\
     &\phantom{50000}\qquad + 33633306)
\\
t(z) &=
	\frac{1}{50000}(357210 z^7 - 1964655 z^6 + 5056506 z^5 - 9737280 z^4
	+ 15174810 z^3\\
     &\phantom{50000}\qquad - 37747395 z^2  + 52924410 z - 21613606),
\end{align*}
wie man mit dem erweiterten euklidischen Algorithmus finden kann.
Bildet man jetzt die Quotienten $s(z)/q(z)$ und $t(z)/q(z)$, heben sich die
polynomiellen Teile weg und es bleibt
\index{polynomieller Teil}%
\[
\frac{r(z)}{q(z)}
=
\frac{2646}{(3z+7)^2}
+
\frac{441z^2-882z+392}{(z-1)^4}.
\]
Der zweite Term hat noch nicht die von der Partialbruchzerlegung verlangte
Form.
Dies wird erreicht, indem man den Zähler wiederholt mit Rest durch $z-1$
dividiert:
\[
\left.
\begin{aligned}
441z^2-882z+392 &= (441z-441) (z-1 ) - 49 \\
441z-441 &= 441 (z-1) + 0
\end{aligned}\right\}
\Rightarrow
441z^2-882z+392 = 441 (z-1)^2 - 49.
\]
Daraus ergeben sich jetzt die Partialbrüche zum Teilnenner $z-1$ als
\[
\frac{t(z)}{(z-1)^4}
=
\frac{441}{(z-1)^2} - \frac{49}{(z-1)^4}.
\]
Insgesamt ist also
\[
\frac{r(z)}{q(z)}
=
\frac{2646}{(3z+7)^2}
+
\frac{441}{(z-1)^2} - \frac{49}{(z-1)^4}
\]
die gesuchte Partialbruchzerlegung.
\end{beispiel}


%
% Integration des rationalen Teils
%
\subsubsection{Integration des rationalen Teils}
Ganz analog der Vorgehensweise in der Methode von Bernoulli im
Abschnitt~\ref{buch:ratint:section:bernoulli} können auch die
Integranden mit Potenzen quadratfreier Polynome $q(z)$ im Nenner
auf den Nenner $q(z)$ reduziert werden.

\begin{satz}[Hermite]
\label{buch:ratint:hermite:satz:reduktion}
\index{Satz!von Hermite}%
Sei $q(z)$ ein quadratfreies Polynom und $p(z)$ ein beliebiges
Polynom von kleinerem Grad $\deg q(z)<\deg p(z)$.
Dann gibt es Polynome $s(z)$ und $t(z)$ Polynome mit
\begin{equation}
s(z)\cdot q(z)
+
t(z)\cdot q'(z)
=
p(z),
\label{buch:ratint:hermite:ggtsz}
\end{equation}
für die zudem
\begin{align}
\int
\frac{p(z)}{q(z)^m}
\,dz
&=
\int
\frac{s(z)+t'(z)/(m-1)}{q(z)^{m-1}}
\,dz
+
\frac{t(z)}{(1-m)q(z)^{m-1}}
\label{buch:rating:hermite:satz:hermite:eqn:int}
\end{align}
wird.
\end{satz}

\begin{proof}
Da $q(z)$ quadratfrei ist sind  $q(z)$ und $q'(z)$ teilerfremd.
Nach dem euklidischen Algorithmus für Polynome gibt es daher
Polynome $s(z)$ und $t(z)$ derart, dass \eqref{buch:ratint:hermite:ggtsz}
gilt.
Einsetzen ins Integral ergibt
\begin{align*}
\int
\frac{p(z)}{q(z)^m}
\,dz
&=
\int
\frac{s(z)\cdot q(z)+t(z)\cdot q'(z)}{q(z)^m}
\,dz
\\
&=
\int
\frac{s(z)}{q(z)^{m-1}}
\,dz
+
\int
t(z)
\cdot
\frac{q'(z)}{q(z)^m}
\,dz
\intertext{Auf das zweite Integral kann partielle Integration
angewendet werden.
Das Polynom $p(z)$  wird abgeleitet während der Bruch
$q'(z)/q(z)^m$ integriert wird.
So entsteht}
&=
\int
\frac{s(z)}{q(z)^{m-1}}
\,dz
+
t(z)
\frac{1}{(1-m)q(z)^{m-1}}
-
\int
\frac{t'(z)}{(1-m)q(z)^{m-1}}
\,dz
\\
&=
\int
\frac{s(z)+t'(z)/(m-1)}{q(z)^{m-1}}
\,dz
+
t(z)
\frac{1}{(1-m)q(z)^{m-1}}.
\end{align*}
Damit ist der Satz bewiesen.
\end{proof}

Durch wiederholte Anwendung des
Satzes~\ref{buch:ratint:hermite:satz:reduktion}
kann man die Potenz des Nenners immer weiter reduzieren, bis
nur noch das Integral
\[
\int
\frac{p(z)}{q(z)}
\,dz
\]
bestimmt werden muss, wobei $p(z)$ und $q(z)$ teilerfremd sind,
$\deg p(z)<\deg q(z)$ und $q(z)$ quadratfrei.
Dieser Teil des Integrationsproblems heisst auch der 
\emph{logarithmische Teil} des Integrals.
\index{logarithmischer Teil}%
Die Methode von Hermite vermeidet eine vollständige Faktorisierung
und ersetzt sie durch eine quadratfreie Faktorisierung.
Sie erreicht aber auch nur eine Reduktion auf quadratfreie Nenner.
Für die vollständig Berechnung der Stammfunktion ist nach der Methode
von Bernoulli weiterhin eine Faktorisierung der quadratfreien Faktoren
sein.
Dies wird aber potentiell dadurch vereinfacht, dass alle Nullstellen
eines quadratfreien Polynoms einfach sind.

\begin{beispiel}
Man finde die Stammfunktion
\begin{align*}
F(z)
&=
\int
\frac{p(z)}{q(z)}\,dz
\\
&=
\int
\frac{
441z^7 + 780z^6 - 2861z^5 + 4085z^4 + 7695z^3 + 3713z^2 - 43253z + 24500
}{
9z^6 + 6z^5 - 65z^4 + 20z^3 + 135z^2 - 154z + 49
}
\,dz
\end{align*}
mit den Polynomen
\begin{align*}
p(z)
&=
441z^7 + 780z^6 - 2861z^5 + 4085z^4 + 7695z^3 + 3713z^2 - 43253z + 24500
\\
q(z)
&=
9z^6 + 6z^5 - 65z^4 + 20z^3 + 135z^2 - 154z + 49.
\end{align*}
Zunächst muss man den Bruch mit Reset dividieren und erhält
\[
F(z)
=
\int 49z + 54\,dz
+
\int
\frac{6651z^4 + 3969z^2 - 37338z + 21854}{q(z)}\,dz,
\]
wobei wir den Zähler im zweiten Integral mit $r(z)$ bezeichnen wollen.
Für die weitere Vereinfachung des zweiten Integrals wird jetzt die
quadratfreie Faktorisierung von $q(z)$ benötigt.
Sie ist
\[
q(z)
=
(z-1)^4(3z+7)^2.
\]
In diesem Fall hat die quadratfreie Faktorisierung eine vollständige
Faktorisierung in Linearfaktoren geliefert.
Damit kann jetzt auch die Partialbruchzerlegung gefunden werden,
wie dies in Beispiel~\ref{buch:ratint:hermite:bsp:partialbruch}
bereits geschehen ist.
Sie ist
\begin{equation}
\frac{r(z)}{q(z)}
=
\frac{294}{(z+\frac{7}{3})^2}
+
\frac{441}{(z-1)^2}
-
\frac{49}{(z-1)^4}.
\label{buch:ratint:hermite:bps:hermite:pf}
\end{equation}
Für die einzelnen Summanden ist jetzt noch die Reduktion nach der Methode
von Hermite von Satz~\ref{buch:ratint:hermite:satz:reduktion} durchzuführen.

Für den ersten Term müssen mit dem euklidischen Algorithmus
$s(z)$ und $t(z)$ gefunden werden derart, dass
\[
s(z)(z+{\textstyle\frac{3}{7}}) + t(z)(z+{\textstyle\frac{3}{7}})
=
s(z)(z+{\textstyle\frac{3}{7}}) + t(z) = 294
\]
gilt, nämlich $s(z)=0$ und $t(z)=294$.
Daraus folgt das Integral für den ersten Term als
\[
\int \frac{294}{(z+\frac37)^2}
=
-\frac{294}{z+{\textstyle\frac37}},
\]
wie man auch von der bekannten Integrationsformel erwartet hätte.

Für den zweiten Term muss man $s(z)$ und $t(z)$ finden derart, dass
\[
s(z)(z-1) + t(z)(z-1)' = s(z)(z-1) + t(z) = 441
\]
gilt.
Der euklidische Algorithmus liefert $s(z)=0$ und $t(z)=441$ und daher
\[
\int \frac{441}{(z-1)^2}\,dz
=
-
\frac{441}{z-1}.
\]
Ebenso funktioniert die Methode für den dritten Term.
Man muss $s(z)$ und $t(z)$ finden derart, dass
\[
s(z)(z-1) + t(z)(z-1)' = s(z)(z-1) + t(z) = -49
\]
gilt.
Es folgt $s(z)=0$ und $t(z)=-49$ und damit 
\[
\int \frac{-49}{(z-1)^4} \,dz
=
\frac{49}{3(z-1)^3}.
\]
Damit ist das gesuchte Integral
\[
F(z)
=
-
\frac{294}{z+\frac{3}{7}}
-
\frac{441}{z-1}
+
\frac{49}{3(z-1)^3}
\]
gefunden.
\end{beispiel}

\begin{beispiel}
Man bestimme das Integral
\[
F(z)
=
\int \frac{p(z)}{q(z)} \,dz
=
\int \frac{z+1}{(z^2+1)^3} \,dz.
\]
Der Nenner liegt bereits in der quadratfreien Faktorisierung mit dem
Faktor $q(z)=z^2+1$ vor, somit hat der Integrand bereits die Form der
Partialbruchzerlegung, es bleibt nur noch, die Reduktion nach Hermite 
mit $m=3$ zu finden.
Dazu werden zunächst $s(z)$ und $t(z)$ bestimmt mit
\[
s(z)q(z) + t(z)q'(z) = z+1,
\]
man findet $s(z)=z+1$ und $t(z)=-\frac12(z^2+z)$.
Einsetzen in die Formel~\eqref{buch:rating:hermite:satz:hermite:eqn:int}
ergibt
\begin{align}
\int
\frac{z+1}{(z+1)^3}
\,dz
&=
\int \frac{2z}{3(z^2+1)^2} \,dz
+
\frac{z^2+z}{4z^4+8z^2+4}.
\label{buch:ratint:hermite:bsp:hermite2:eqn:i3}
\end{align}
Der Prozess muss jetzt für das verbleibende Integral auf der rechten
Seite wiederholt werden.
Diesmal findet man
\[
s(z)=\frac{2z+3}{4}
\qquad\text{ und }\qquad
t(z)=-\frac{2z^2+3}{8}
\]
und nach Einsetzen
\begin{align}
\int \frac{2z}{3(z^2+1)^2} \,dz
&=
\int \frac{\frac{3}{8}}{z^2+1}\,dz
+
\frac{2z^2+3z}{8z^2+8}.
\label{buch:ratint:hermite:bsp:hermite2:eqn:i2}
\end{align}
Kombiniert man
\eqref{buch:ratint:hermite:bsp:hermite2:eqn:i3}
und
\eqref{buch:ratint:hermite:bsp:hermite2:eqn:i2}
findet man
\begin{align*}
\int
\frac{z+1}{(z+1)^3}\,dz
&=
\frac38 \int \frac{1}{z^2+1}\,dz
+
\frac{2z^4+3z^3+4z^2+5z}{8z^4+16z^2+8}
\\
&=
\arctan z
+ 
\frac{3z^3+5z-2}{8z^4+16z^2+8}
+\frac14.
\end{align*}
Der Term $\frac14$ auf der echten Seite ist entstanden durch Division
des Zählers mit Rest durch den Nenner.
Er kann in der Integrationskonstanten absorbiert werden.
\end{beispiel}
